\chapter{Аналитическая часть}

\section{Рекурсия}

Согласно общепринятому определению рекурсии~\cite{sweygart}, рекурсией называют способ задания объекта, при котором он определяется через самого себя. В программировании рекурсивной является функция, вызывающая саму себя для решения подзадачи, аналогичной основной. Такой подход позволяет естественным образом описывать процессы, обладающие свойством самоподобия.

Формально рекурсивная функция может быть записана следующим образом~(\ref{eq:rec_def}):
\begin{equation}
	F(x)=
	\begin{cases}
		F_0(x), & x \in B, \\[3pt]
		G\left(x, F(h(x))\right), & x \notin B,
	\end{cases}
	\label{eq:rec_def}
\end{equation}
где $B$ --- множество базовых случаев, при которых дальнейшее разбиение задачи не требуется, $h(x)$ --- преобразование аргумента, ведущее к подзадаче, а $G$ --- правило получения результата для текущего уровня вычисления.

Корректная рекурсивная конструкция всегда включает:
\begin{itemize}
	\item по крайней мере один базовый случай;
	\item механизм гарантированного приближения к базовому случаю;
\end{itemize}

\section{Графы как математическая модель алгоритмов}

Математически граф $G$ определяется множеством вершин~(\ref{eq:graph_vertices}):
\begin{equation}
	X = \{ x_1, x_2, \ldots, x_n \},
	\label{eq:graph_vertices}
\end{equation}

и множеством рёбер~(\ref{eq:graph_edges}):
\begin{equation}
	A = \{ a_1, a_2, \ldots, a_m \}.
	\label{eq:graph_edges}
\end{equation}

Таким образом, граф задаётся парой множеств~(\ref{eq:graph_def}):
\begin{equation}
	G = (X, A).
	\label{eq:graph_def}
\end{equation}

Использование графов для анализа алгоритмов позволяет:
\begin{itemize}
	\item описывать структуру вычислений;
	\item фиксировать зависимости по управлению;
	\item выявлять зависимости по данным;
	\item определять независимые компоненты алгоритма;
	\item находить возможные области для параллельного выполнения;
\end{itemize}

\section{Графовые модели алгоритмов}

В данной работе применяются четыре графовые модели: граф управления, информационный граф, операционная история и информационная история;

\subsection{Граф управления}

Граф управления (Control Flow Graph, CFG) --- ориентированный граф~(\ref{eq:cfg_def}):
\begin{equation}
	G_{\text{ctrl}} = (V_{\text{ctrl}}, A_{\text{ctrl}}),
	\label{eq:cfg_def}
\end{equation}
где:
\begin{itemize}
	\item $V_{\text{ctrl}}$ --- множество операторов программы;
	\item $A_{\text{ctrl}}$ --- множество дуг, отражающих возможные переходы управления;
\end{itemize}

Дуги показывают последовательность выполнения операций, а также ветвления, циклы и рекурсивные возвраты. Если между двумя вершинами отсутствует путь, эти части алгоритма независимы по управлению.

\subsection{Информационный граф}

Информационный граф (Data Flow Graph, DFG) --- ориентированный граф~(\ref{eq:dfg_def}):
\begin{equation}
	G_{\text{data}} = (V_{\text{ops}}, A_{\text{data}}),
	\label{eq:dfg_def}
\end{equation}
где вершины соответствуют операциям, а дуга $(u, v)$ существует, если результат операции $u$ используется операцией $v$.

Такой граф показывает зависимости по данным. Отсутствие зависимости между двумя операциями означает возможность их параллельного выполнения.

\subsection{Операционная история}

Операционная история представляет собой частично упорядоченное множество~(\ref{eq:ops_hist}):
\begin{equation}
	H_{\text{ops}} = (O, \prec),
	\label{eq:ops_hist}
\end{equation}
где $O$ --- множество операций, а $\prec$ --- отношение строгого частичного порядка, определяющее их обязательный порядок выполнения.

Если для двух операций отношение порядка не задано, их можно выполнять независимо или параллельно.

\subsection{Информационная история}

Информационная история описывает изменения данных и задаётся структурой~(\ref{eq:data_hist}):
\begin{equation}
	H_{\text{data}} = \{ (d, t_{\text{create}}, U(d)) \},
	\label{eq:data_hist}
\end{equation}
где:
\begin{itemize}
	\item $d$ --- значение данных или переменная;
	\item $t_{\text{create}}$ --- операция, породившая это значение;
	\item $U(d)$ --- множество операций, использующих его;
\end{itemize}

Эта модель позволяет выявлять конфликты при обращении к данным, определять независимые цепочки преобразований и устанавливать возможность их параллельного выполнения.

\section*{Вывод}

В данном разделе были рассмотрены понятия рекурсии и графов.

\clearpage
